% Purpose: Conference abstract on education, artificial intelligence, and
% future leadership.
% Status: Revised Draft; author and venue review pending.
% Source-of-truth: Author-provided topic, pedagogical framing, and
% philosophical debate.
\documentclass[12pt]{article}
\usepackage{amsmath}
\usepackage{microtype}

\title{Education in the Age of AI:\break\large Cultivating Future Thinkers (?)}
\author{Tanaka Hernandez, Ken}
\date{\today}

\begin{document}
\maketitle

\begin{abstract}
  Education has repeatedly confronted ``\textit{technological advancements}''
  that have changed how knowledge is recorded, accessed, and shared: from
  ancient cave paintings, to \textit{Phaedrus}' passage warning against
  standardized (Gutenberg press)-like imprints following the industrial
  Prussian mass-schooling, to the dispute over electronic calculators in
  mathematics, and the distributed digital memory of the internet. Artificial
  Intelligence (AI)'s \textit{two-sided pseudo-cognitive} prowess, re-opens the
  debate, asking \textbf{what} composes human intellect and \textbf{how} can it
  be incorporated to strengthen learning and prepare future leaders.

  This lecture avoids any preference over pedagogical doctrine, instead
  presents comparative insights on how to \textbf{harness} modern AI tools as
  potential support, by delegating mechanical friction in exploration,
  feedback, and routine tasks; to allow conscious reallocation of cognition
  toward \textbf{\textit{Conatus ad Novandum}}: the deliberate strive to
  question, create, and innovate that defines \textbf{humane} visionary
  leadership.

  The lecture concludes with a biological and epistemological challenge
  (\textit{a priori} vs. \textit{a posteriori}): humanity has \textit{evolved}
  extensive cognitive capabilities (\textbf{\textit{cogito}}), yet remains
  within sensory, embodied, and affective experience
  (\textbf{\textit{sentio}}). Can a purely disembodied ``cognitive animal''
  exist, or is it merely a mathematical artifact, and if not, why have we been
  following symbolic abstractions to enable such environmental manipulation for
  our self convenient \textit{greed}? \textit{Evolution} (from a
  non-deterministic, non-teleological perspective) suggests an evolutionary
  \textit{split}? Which follows the question: what is \textbf{\textit{life}}?
\end{abstract}

\noindent\textbf{Keywords:} artificial intelligence; agentic harness;
educational leadership; historical pedagogies; cognitive offloading; conatus ad
novandum; embodied cognition; cogito~vs.~sentio
\newpage
\begin{abstract}
  La educación se ha enfrentado reiteradamente a los ``\textit{avances
  tecnológicos}'' que han transformado la manera de registrar, acceder y
  compartir el conocimiento: desde las pinturas rupestres, seguido de la
  advertencia en el \textit{Fedro} contra las impresiones estandarizadas
  (semejantes a las de la imprenta de Gutenberg), siguiendo el modelo
  industrial prusiano de escolarización masiva, hasta la controversia sobre las
  calculadoras electrónicas en matemáticas y la memoria digital distribuida de
  internet. La capacidad \textit{pseudo-cognitiva dual} de la Inteligencia
  Artificial (IA) reabre el debate y plantea \textbf{qué} compone el intelecto
  humano y \textbf{cómo} puede incorporarse para fortalecer el aprendizaje y
  preparar a los futuros líderes.

  Esta conferencia no privilegia ninguna doctrina pedagógica en cambio,
  presenta comparativas sobre cómo implementar \textbf{arneses} a las
  herramientas de IA como posible apoyo, delegándoles la fricción mecánica de
  la exploración, retroalimentación y tareas rutinarias; permitiendo así,
  reasignar conscientemente los recursos cognitivos hacia
  \textbf{\textit{Conatus ad Novandum}}: el esfuerzo deliberado por cuestionar,
  crear e innovar que define un liderazgo visionario \textbf{humano}.

  La conferencia concluye con un desafío biológico y epistemológico
  (\textit{a priori} vs. \textit{a posteriori}): la humanidad ha
  \textit{evolucionado} amplias capacidades cognitivas
  (\textbf{\textit{cogito}}), pero sigue inmersa en una experiencia sensorial,
  encarnada y afectiva (\textbf{\textit{sentio}}). ¿Puede existir un animal
  puramente ``cognitivo'' e incorpóreo, o es tan solo un artefacto matemático
  y, de no ser así, ¿por qué hemos recurrido a abstracciones simbólicas para
  manipular el entorno en beneficio de nuestra \textit{codicia}? ¿Acaso la
  \textit{evolución} (desde una perspectiva no determinista ni teleológica)
  sugiere una posible \textit{bifurcación} evolutiva? En el cual sigue la
  pregunta: ¿qué es la \textbf{\textit{vida}}?
\end{abstract}

\noindent\textbf{Keywords:} inteligencia artificial; arnés agéntico; liderazgo
educativo; pedagogías históricas; descarga cognitiva; conatus ad novandum;
cognición encarnada; cogito~vs.~sentio

\end{document}
